\chapter{Exchange Symmetry, Spin, and the Two Minus Signs}

\section{Shell Filling and the Spin--Statistics Question}

The shell-model review begins with two properties associated with electrons. An electron has spin one-half, with the two dimensionless spin projections
\begin{equation}
m_s=+\frac12,
\qquad
m_s=-\frac12,
\end{equation}
and electrons obey the Pauli exclusion principle. These facts are related, but they are not the same statement.

The phrase ``the same state'' needs care. No two electrons may occupy the same complete one-particle state. Two electrons may nevertheless share one spatial orbital when their spin states differ. In the independent-electron approximation to helium, both electrons occupy the lowest spatial orbital, one with spin up and the other with spin down. More precisely, the alternatives ``the first is up and the second is down'' and ``the first is down and the second is up'' belong to an entangled two-electron spin state.

This is only the old shell-model approximation. It ignores the interaction between the electrons and is being used for its simple counting. For lithium, the first two electrons fill the lowest spatial orbital with opposite spin. The third cannot enter either of those complete states and must occupy a higher orbital.

The next step is the familiar orbital-angular-momentum diagram. The sectors with
\begin{equation}
\ell=0,\qquad \ell=1,\qquad \ell=2
\end{equation}
are called \(S\), \(P\), and \(D\) states. A fixed value of \(\ell\) has
\begin{equation}
2\ell+1
\end{equation}
orbital states, so
\begin{equation}
\ell=0\Rightarrow 1,
\qquad
\ell=1\Rightarrow 3,
\qquad
\ell=2\Rightarrow 5.
\end{equation}
The Coulomb problem has additional degeneracies: radial levels from different \(\ell\)-sectors can have the same energy. For the shell-filling argument, the next shell contains one \(S\) orbital and three \(P\) orbitals. It therefore has
\begin{equation}
1+3=4
\end{equation}
spatial states. Spin doubles the count:
\begin{equation}
2(1+3)=8.
\end{equation}

Those eight complete one-electron states are filled in the simple progression from lithium through neon. Beyond the closed shell, sodium again has one outer electron, which is why this crude picture makes it resemble lithium. The chemistry is not the destination. Its point is that half-integer spin and exclusion appear together.

Photons provide the contrast. They have integer spin and do not obey an exclusion rule. A classical electromagnetic wave corresponds to many photons occupying the same mode. Bosons can accumulate in a common state, whereas identical fermions cannot occupy one complete state more than once. The question is whether those two pairings could have been otherwise: half-integer spin without exclusion, or integer spin with exclusion. In three spatial dimensions the spin--statistics relation says no. The familiar statement requires qualification in one and two spatial dimensions, and no proof of the field-theoretic result is attempted here.

\begin{classroomqa}
\lecturetimestamp{00:11:48}
\audiencequestion{Is there a maximum number of bosonic quanta that can occupy one state?}
\lecturerresponse{No. Bosons have no occupancy limit, and they tend to accumulate in the same state.}
\end{classroomqa}

\section{Many-Particle Wavefunctions}

The connection will be approached through two minus signs, one from exchange and one from rotation. Begin with exchange, and temporarily suppress spin.

Let \(x_i\) denote the position of particle \(i\). The subscript is a particle label, not a Cartesian component; the three spatial coordinates are compressed into the single symbol \(x_i\). A position basis for \(n\) particles is
\begin{equation}
\lvert x_1,x_2,\ldots,x_n\rangle.
\end{equation}
For an abstract state \(\lvert\Psi\rangle\), the coordinate wavefunction is
\begin{equation}
\psi(x_1,x_2,\ldots,x_n)
=
\langle x_1,x_2,\ldots,x_n\mid\Psi\rangle.
\end{equation}
Its squared magnitude,
\begin{equation}
\bigl\lvert\psi(x_1,x_2,\ldots,x_n)\bigr\rvert^2,
\end{equation}
is the joint probability density for finding particle \(1\) at \(x_1\), particle \(2\) at \(x_2\), and so on.

Now suppose the particles are identical. Imagine first that one electron is painted blue and another red. Then ``blue at \(x_1\), red at \(x_2\)'' could be distinguished from ``blue at \(x_2\), red at \(x_1\).'' Actual electrons carry no such colors or individual tags. Interchanging the bookkeeping labels cannot create a new physical arrangement.

An overall phase does not change probabilities or expectation values. For a state in a definite exchange sector, the coordinate wavefunction may therefore obey
\begin{equation}
\psi(x_2,x_1)=\eta\,\psi(x_1,x_2),
\qquad
\lvert\eta\rvert=1.
\end{equation}
Applying the same exchange twice restores the original arguments and gives
\begin{equation}
\eta^2=1,
\end{equation}
so the two possible exchange eigenvalues are
\begin{equation}
\eta=+1
\qquad\hbox{or}\qquad
\eta=-1.
\end{equation}

\editorialnote{The double-exchange algebra determines the possible eigenvalues only after a state has been taken to have definite exchange behavior. The further restriction that the allowed states of a fixed species lie in one exchange sector is a physical rule; it does not follow from \(\eta^2=1\) alone. The same qualification applies to the exclusion argument below: antisymmetry implies exclusion once the fermionic sector has been selected, but the algebra of exchange by itself does not select that sector.}

For bosons the allowed wavefunctions are symmetric,
\begin{equation}
\psi_B(x_1,x_2)=\psi_B(x_2,x_1),
\end{equation}
whereas for fermions they are antisymmetric,
\begin{equation}
\psi_F(x_1,x_2)=-\psi_F(x_2,x_1).
\end{equation}
The minus sign is invisible when it multiplies the whole state, but it matters when amplitudes are added. That is how exchange symmetry acquires observable consequences.

\section{Products, Sums, and Exclusion}

Consider first two particles assigned to the same one-particle orbital \(\psi_0\). Their product is
\begin{equation}
\psi_0(x_1)\psi_0(x_2).
\end{equation}
Exchanging the particle coordinates gives
\begin{equation}
\psi_0(x_2)\psi_0(x_1)
=
\psi_0(x_1)\psi_0(x_2).
\end{equation}
Wavefunction values are numbers, not operators, so the factors commute. The same-state product is symmetric. It is an allowed bosonic form, but it cannot describe two identical fermions.

Next assign the particles to different one-particle orbitals:
\begin{equation}
\psi_0(x_1)\psi_1(x_2).
\end{equation}
Exchange produces
\begin{equation}
\psi_0(x_2)\psi_1(x_1).
\end{equation}
For distinct functions the exchanged product is generally neither the original expression nor its negative. The bare product therefore has neither required exchange symmetry. This does not forbid one particle from occupying \(\psi_0\) while another occupies \(\psi_1\). It means that the two-particle state must combine the two assignments.

The symmetric combination is
\begin{equation}
\psi_B^{(01)}(x_1,x_2)
\propto
\psi_0(x_1)\psi_1(x_2)
+
\psi_0(x_2)\psi_1(x_1),
\end{equation}
and the antisymmetric combination is
\begin{equation}
\psi_F^{(01)}(x_1,x_2)
\propto
\psi_0(x_1)\psi_1(x_2)
-
\psi_0(x_2)\psi_1(x_1).
\end{equation}
Proportionality signs are sufficient because no normalization is needed for the exchange argument. Both states contain one particle in each orbital; they differ only in how the two indistinguishable assignments are combined.

Now let the orbitals coincide. The bosonic combination becomes
\begin{align}
\psi_B^{(00)}(x_1,x_2)
&\propto
\psi_0(x_1)\psi_0(x_2)
+
\psi_0(x_2)\psi_0(x_1)\\
&=
2\psi_0(x_1)\psi_0(x_2),
\end{align}
which is nonzero. The fermionic combination gives
\begin{align}
\psi_F^{(00)}(x_1,x_2)
&\propto
\psi_0(x_1)\psi_0(x_2)
-
\psi_0(x_2)\psi_0(x_1)\\
&=0.
\end{align}
With the antisymmetric exchange rule, exclusion follows: there is no nonzero antisymmetric state with both fermions in the same one-particle state.

For a many-particle state the rule applies to any transposition of two arguments. A transposition leaves a bosonic wavefunction unchanged and reverses the sign of a fermionic one. Repeating the same transposition is trivial in both cases. The surprising part is that a single exchange, though it changes no observable particle labels, changes the sign of a fermionic state.

\section{Spin in the Complete Particle Label}

Spin can now be restored with a small change. A one-particle amplitude depends on both position and a discrete spin label,
\begin{equation}
\psi(x,\sigma),
\end{equation}
where \(\sigma=+1\) or \(-1\) labels the chosen \(\sigma_z\) result. The squared magnitude is the joint probability density for the stated position and spin label.

For many particles the wavefunction is
\begin{equation}
\psi(x_1,\sigma_1;\,x_2,\sigma_2;\,\ldots).
\end{equation}
Exchange swaps everything belonging to the two particles, not only their positions:
\begin{align}
\psi_B(x_1,\sigma_1;\,x_2,\sigma_2)
&=
+\psi_B(x_2,\sigma_2;\,x_1,\sigma_1),\\
\psi_F(x_1,\sigma_1;\,x_2,\sigma_2)
&=
-\psi_F(x_2,\sigma_2;\,x_1,\sigma_1).
\end{align}
This is the precise form of the opening shell-model statement. Two electrons may share one spatial orbital only when their complete states, including spin, are different.

\begin{classroomqa}
\lecturetimestamp{00:34:33}
\audiencequestion{Is exchange symmetry an instance of the permutation group?}
\lecturerresponse{Yes. Exchanging particles is a permutation, and the plus or minus transformation law gives a representation of the permutation group.}
\end{classroomqa}

\begin{classroomqa}
\lecturetimestamp{00:35:12}
\audiencequestion{Why is the two-particle construction a product rather than a sum of the one-particle wavefunctions?}
\lecturerresponse{A sum describes one particle in a superposition. A product belongs to the two-particle tensor-product space and describes one particle assigned to each one-particle state.}
\end{classroomqa}

The full-label exchange is shown on the board below.\footnote{The ordering of the position--spin labels is reconstructed from the legible blackboard frame and the surrounding explanation from 00:33:43 to 00:34:25.}
\begin{figure}[htbp]
\centering
\includegraphics[width=0.9\linewidth]{lecture_05_frame_03.png}
\caption{Exchanging two identical particles swaps their complete position--spin labels. A bosonic state is unchanged, while a fermionic state changes sign.}
\lectureframe{00:35:25}
\end{figure}

These are different uses of linearity. For one particle,
\begin{equation}
\alpha\psi_0(x)+\beta\psi_1(x)
\end{equation}
is a superposition of two alternatives. For two particles,
\begin{equation}
\psi_0(x_1)\psi_1(x_2)
\end{equation}
is a product amplitude on a space with two particle coordinates. Symmetrization or antisymmetrization combines that product with its exchanged partner.

\begin{classroomqa}
\lecturetimestamp{00:36:13}
\audiencequestion{Does exchanging particles produce an interference effect?}
\lecturerresponse{Yes. The vanishing same-state fermionic amplitude can be understood as destructive interference between the two exchanged contributions.}
\end{classroomqa}

\begin{classroomqa}
\lecturetimestamp{00:37:13}
\audiencequestion{Did Pauli derive exclusion using this same symmetry argument?}
\lecturerresponse{Pauli began from an empirical fact and found a mathematical framework in which it could be expressed. He was not initially deriving exclusion from a more fundamental theory.}
\end{classroomqa}

Antisymmetry supplies the mathematical framework in which the empirical exclusion rule is expressed.

\begin{classroomqa}
\lecturetimestamp{00:39:01}
\audiencequestion{Is particle-exchange interference connected with interference between alternative paths?}
\lecturerresponse{Yes. One history has the two particles pass without crossing, while another has them cross into the exchanged destinations. The amplitudes add for bosons and can cancel for fermions.}
\end{classroomqa}

Exchange supplies the first minus sign. Rotation supplies the second.

\section{The Second Minus Sign: Rotation by \(2\pi\)}

A full rotation appears to do nothing. Every ordinary geometric direction returns to its starting point after an angle \(2\pi\). What must return, however, are the probabilities and expectation values. The state vector itself may return with an overall phase, just as it may under exchange.

\begin{classroomqa}
\lecturetimestamp{00:42:12}
\audiencequestion{Which axis is used for the rotation, and is the generator orbital angular momentum, spin, or both?}
\lecturerresponse{One may choose an external axis; below it is taken to be the \(z\)-axis. The generator is total angular momentum, including both orbital and spin contributions. The spin itself need not point along the chosen axis.}
\end{classroomqa}

\editorialnote{The angular-momentum formulas in this discussion use \(\hbar=1\), so \(m\) is the dimensionless eigenvalue of \(J_z\).}

Let \(\psi(\theta)\) denote the state obtained as the rotation parameter advances through an angle \(\theta\). In the convention used here,
\begin{equation}
J_z\psi(\theta)
=
-i\frac{\partial}{\partial\theta}\psi(\theta).
\end{equation}
For a state with definite \(J_z\),
\begin{equation}
J_z\psi(\theta)=m\psi(\theta).
\end{equation}
Combining the equations gives
\begin{equation}
-i\frac{\partial}{\partial\theta}\psi(\theta)
=
m\psi(\theta),
\end{equation}
with solution
\begin{equation}
\psi(\theta)=e^{im\theta}\psi(0).
\end{equation}

The allowed values of \(m\) are integers or half-integers. Therefore
\begin{equation}
\psi(2\pi)
=
\begin{cases}
+\psi(0), & m\in\mathbb Z,\\[4pt]
-\psi(0), & m\in\mathbb Z+\tfrac12.
\end{cases}
\end{equation}
After another full turn,
\begin{equation}
\psi(4\pi)=\psi(0)
\end{equation}
in either case. This is the second minus sign: a state in a half-integer angular-momentum multiplet changes sign under a \(2\pi\) rotation.

The sign does not alter \(\lvert\psi\rvert^2\), so it is not visible when it multiplies the entire state. Even so, a \(2\pi\) rotation is not represented by the identity on a half-integer-spin state. The distinction becomes observable when the sign is made relative to another coherent amplitude.

\section{The \(2\pi\) and \(4\pi\) Paths}

The distinction between \(2\pi\) and \(4\pi\) can be made visible without further algebra. Hold a cup in one hand and turn it through one full revolution without releasing it. The cup returns to its original orientation, but the arm is twisted. Continue through a second revolution and the arm can be returned to its original position without reversing the cup's motion.

A related construction places a sphere inside a box and joins it to the box with strings. Turning the sphere through \(2\pi\) leaves a tangle. In ordinary three-dimensional space, the tangle associated with a \(4\pi\) turn can be unwound. The relevant object is not merely the final orientation of the sphere but the path followed through the space of rotations.

\begin{classroomqa}
\lecturetimestamp{00:54:22}
\audiencequestion{If the sphere is turned twice about the same axis, can the strings really be untangled? Does that require another spatial dimension?}
\lecturerresponse{The \(4\pi\) path can be unwound in ordinary three-dimensional space. Time parametrizes the motion; it is not another spatial direction. A \(2\pi\) path cannot be unwound in this way, while a \(4\pi\) path can.}
\end{classroomqa}

The demonstrations distinguish a nontrivial \(2\pi\) path from an unwindable \(4\pi\) path without developing a formal theory of the rotation group.

\section{Exchange and Rotation in the Belt Picture}

The two minus signs now stand side by side. Exchanging two identical fermions changes the sign of their joint wavefunction. Rotating one half-integer-spin state through \(2\pi\) also changes its sign. A belt gives a geometrical relation between the two operations.

Arrange the belt so that its two lower ends have been exchanged. Viewed another way, the same belt carries a \(2\pi\) twist. Performing the exchange together with the \(2\pi\) rotation removes the twist and leaves an untwisted belt. The combined process can be continuously reduced to doing nothing.\footnote{The exchange--rotation observation is attributed to David Finkelstein in the lecture. No specific publication is identified in the source record, so the name is retained here only as that attribution.}

\begin{classroomqa}
\lecturetimestamp{01:01:12}
\audiencequestion{Is this only a statement about how the belt is embedded in space rather than a statement about particle topology?}
\lecturerresponse{The belt can be read as a particle history in spacetime. With time running upward, a cross-section through it gives particle-like objects. Pairs are created, one object is rotated through \(2\pi\), and two identical objects are exchanged. In that history the rotation and exchange phases cancel.}
\end{classroomqa}

The history makes the phase bookkeeping explicit. Begin with no particles. At two separated places, create a particle--antiparticle pair, leaving four particles at an intermediate time. Rotate one particle through \(2\pi\) while exchanging two identical particles elsewhere in the configuration. The complete history can be deformed into the elementary history in which a pair is created and later annihilated. In the sign bookkeeping under discussion, that trivial history has no additional phase. If \(\eta_{\mathrm{rot}}\) and \(\eta_{\mathrm{ex}}\) denote the two signs, the spoken cancellation can be written compactly as\footnote{This phase equation records the cancellation described from 01:06:50 to 01:07:18; the lecture presents the construction as suggestive rather than as a rigorous proof.}
\begin{equation}
\eta_{\mathrm{rot}}\eta_{\mathrm{ex}}=1.
\end{equation}

\begin{classroomqa}
\lecturetimestamp{01:03:34}
\audiencequestion{Can exchange be related to rotations by letting each of two particles move through \(\pi\) and then reorienting them?}
\lecturerresponse{That construction may reflect related mathematics, but it uses two separate \(\pi\) rotations. The process being analyzed here contains one \(2\pi\) rotation, so the two arguments are not identical.}
\end{classroomqa}

The argument depends on the possibility of creating and annihilating particle-like objects. It would not have this form in a theory whose particles were eternal. It is also not offered as the standard proof of the spin--statistics theorem. Its purpose is to display a topological reason that the exchange and rotation signs fit together.

\section{Lumps, Mirror Images, and Annihilation}

The construction can be repeated with localized field configurations. Suppose one or more fields take their vacuum values almost everywhere but have a nonzero shape and structure in a small region. That localized configuration is a lump.

Now make a second configuration as the appropriate mirror image of the first. Call one the particle and the other the antiparticle. Bring them together. Corresponding pieces of the field configurations meet and undo one another. At an intermediate stage some parts of the lumps may already have disappeared while other parts remain. Continuing the deformation removes the remaining structure and returns the fields to the vacuum configuration. Two oppositely arranged twists in a belt give the same picture: push the twists together and they unwind.

The process can be run backward. Empty space can produce a lump and its mirror partner. One may create two such pairs, rotate one lump through \(2\pi\), exchange two identical lumps, and annihilate the pairs again. The same exchange--rotation cancellation then applies. This remains a heuristic extension of the construction, not a theorem about every localized excitation.

\begin{classroomqa}
\lecturetimestamp{01:11:30}
\audiencequestion{Are there physical realizations of these lump-like field configurations?}
\lecturerresponse{Yes. Examples include nuclei in Skyrmion-type descriptions, magnetic monopoles, and other particle-like excitations of fields.}
\end{classroomqa}

The mirrored-lump picture is what makes the generalization intelligible. The relevant feature is not that the object began as an elementary point particle, but that a particle-like excitation and its partner can be created, brought together, and annihilated continuously. That feature motivates a broader relation between rotational and exchange behavior.

\begin{classroomqa}
\lecturetimestamp{01:16:12}
\audiencequestion{Does spin one-half mean that the state accumulates only half the usual rotation phase?}
\lecturerresponse{The phase in this convention is \(e^{im\theta}\). For half-integer \(m\), a \(2\pi\) rotation gives \(-1\), while a \(4\pi\) rotation gives \(+1\).}
\end{classroomqa}

\section{A Global Sign and a Null Experiment}

Can the minus sign from a \(2\pi\) rotation be observed? Begin with an experiment in which it cannot. Trap an electron in a magnetic field, let its spin follow the field while the apparatus is turned slowly through \(2\pi\), and then release the electron into an interference apparatus. Compare this with the same procedure without the rotation.

The rotation changes the state by
\begin{equation}
\psi\longmapsto-\psi,
\end{equation}
but the minus sign multiplies the whole state. Every probability is unchanged, so the final interference pattern retains no record of whether the preliminary rotation occurred. This is a null experiment: a common phase has nothing with which to interfere.

\begin{classroomqa}
\lecturetimestamp{01:20:01}
\audiencequestion{Why must the magnetic-field rotation be performed slowly?}
\lecturerresponse{The spin must continue to follow the rotating magnetic field. If the field is dragged too quickly, the spin will not remain locked to it.}
\end{classroomqa}

\section{One Electron in Two Coherent Branches}

The sign becomes observable when it is made relative. Send one electron through a beam splitter. Part of its wavefunction follows one route and part follows another:
\begin{equation}
\psi(x)=\psi_1(x)+\psi_2(x).
\end{equation}
This is one electron, not two electrons. It is a coherent superposition of two alternatives. The distinction must be kept clear: the two packets are not two independent particles.

Guide the two packets into separate regions containing magnetic fields. Rotate the field in one region through \(2\pi\), leaving the other unchanged. For a half-integer-spin state,
\begin{equation}
\psi_1(x)+\psi_2(x)
\longmapsto
\psi_1(x)-\psi_2(x),
\end{equation}
where the label of the rotated branch has been chosen as \(2\). The sign is now relative. When the two packets later overlap, addition has been changed into subtraction. A point of constructive interference can become a point of destructive interference, or even a node.

\begin{classroomqa}
\lecturetimestamp{01:23:46}
\audiencequestion{What is being rotated if there is only one electron but two wavefunction branches?}
\lecturerresponse{The electron is in a coherent superposition of two localized alternatives. One branch is rotated and the other is not. If the boxes are opened to learn which alternative occurred, that measurement destroys the interference.}
\end{classroomqa}

Opening the boxes would reveal only one electron, in one box or the other. Without that measurement, both amplitudes remain in the state. Repeating the experiment builds up the interference pattern one detection at a time. At a node produced by the relative minus sign, no electron is detected.

\section{Magnetic Precession and Two-Slit Interference}

A more practical schematic arrangement replaces the very slow box experiment. Split a coherent electron wave into two paths. In one path, let the electron pass through a magnetic region with its spin initially transverse to the field. The spin then precesses. By choosing the field strength and transit time, the spin in that branch can execute one full \(2\pi\) turn, while the other branch undergoes no such turn.\footnote{The lecture states that a related spinor-phase experiment had been performed and describes its principle, but the source record does not identify a paper or a particular apparatus. The account here is therefore limited to the construction given in the lecture.}

\begin{classroomqa}
\lecturetimestamp{01:27:05}
\audiencequestion{How can the path of the single electron be determined?}
\lecturerresponse{It cannot be determined without ruining the experiment. Any measurement that records which path the electron took destroys the interference.}
\end{classroomqa}

\begin{classroomqa}
\lecturetimestamp{01:27:50}
\audiencequestion{What exactly is a branch of the wavefunction, and how does the splitter produce two branches?}
\lecturerresponse{A branch is a wave packet localized along one route. The splitter turns the incoming wave into two spatially separated packets, and electric or magnetic fields can guide those packets along different paths.}
\end{classroomqa}

Apart from the magnetic region, this is the two-slit experiment. An incoming electron wave separates into two packets. The outgoing waves spread into a common detection region, where their amplitudes overlap. With no relative rotation the central amplitudes may add. A \(2\pi\) spin rotation in one route reverses one amplitude and can make the same contributions cancel.

\begin{classroomqa}
\lecturetimestamp{01:30:46}
\audiencequestion{How can there be interference if the electron travels through distinct tubes or slits?}
\lecturerresponse{The final detection does not reveal the path. The two outgoing wave amplitudes spread and overlap at the detector. Looking to see which path was taken would destroy the interference.}
\end{classroomqa}

\begin{classroomqa}
\lecturetimestamp{01:33:22}
\audiencequestion{Must the tube or slit diameter be comparable to the electron wavelength?}
\lecturerresponse{For a literal diffraction-based two-slit realization, the opening must be sufficiently small to produce the needed spreading. The apparatus being described is schematic rather than a practical design.}
\end{classroomqa}

\begin{classroomqa}
\lecturetimestamp{01:33:44}
\audiencequestion{Would a corresponding \(2\pi\) polarization rotation of photons leave the interference pattern unchanged?}
\lecturerresponse{Yes. Photons have integer spin, so a full \(2\pi\) rotation does not introduce the fermionic minus sign.}
\end{classroomqa}

\begin{classroomqa}
\lecturetimestamp{01:34:12}
\audiencequestion{What is the slight departure of the practical experiment from the idealized version?}
\lecturerresponse{The ideal construction rotates the spin adiabatically. In the practical description the electron moves through a magnetic region and its spin precesses more quickly, but the measured phase is still the phase associated with a \(2\pi\) turn.}
\end{classroomqa}

The slow construction isolates the idea by ensuring that the spin follows the field without additional motion. The precession construction reaches the same relative sign dynamically. In either version, the two routes must cease to be distinguishable before an interference pattern can appear.

\begin{classroomqa}
\lecturetimestamp{01:36:15}
\audiencequestion{Must the two branches be recombined after the splitter?}
\lecturerresponse{Yes. The emerging wave packets must overlap, or spread onto a common detection region, so that their amplitudes can form an interference pattern.}
\end{classroomqa}

\begin{classroomqa}
\lecturetimestamp{01:38:12}
\audiencequestion{Is this one-electron rotation experiment a consequence of Pauli exclusion?}
\lecturerresponse{No. There is only one electron. The experiment measures the spinor phase under rotation, not the exclusion of two fermions from one state.}
\end{classroomqa}

\section{Exchange Interference and Relative Rotation}

An exchange sign could be exposed by the same general method, although the proposed apparatus is more complicated. Begin with two electrons and split their joint wavefunction into two coherent branches. There are still only two electrons, not four: the pair follows one alternative or the other. Exchange the two electrons in one branch and leave them unexchanged in the other. After the branches are brought back together, the fermionic exchange sign appears as a relative phase and can change the interference.

The final proposal compares the two minus signs directly. One branch acquires the fermionic exchange phase, while the other acquires the corresponding \(2\pi\)-rotation phase. Both phases are \(-1\), so the relative effect cancels.\footnote{The recording between 01:40:26 and 01:41:15 clearly gives an exchange phase, a \(2\pi\)-rotation phase, and their cancellation, but the detailed assignment of the rotation is not recoverable from the available speech.} The lecture presents this as a proposed thought experiment and says that, to the speaker's knowledge, it had not been carried out.

\begin{classroomqa}
\lecturetimestamp{01:41:34}
\audiencequestion{Would observing an exchange phase challenge the indistinguishability of the two electrons?}
\lecturerresponse{No. Quantum indistinguishability is subtler than classical sameness. It permits the state vector to change sign under exchange even though no observable label distinguishes the particles.}
\end{classroomqa}

There is one last qualification. Rotating an entire isolated world through \(2\pi\) would give only a common phase, and no internal experiment could reveal it. The detectable experiment rotates one coherent branch relative to another. What is observed is therefore not an absolute sign but the interference produced by a relative sign.

